\section{Discussion}\label{sec:discussion}
\subsection{When a minimal, agentic loop is sufficient}
The reference deployment suggests three conditions under which \ffrs{} as described is enough. First, the project already lives in one tracker: routing is then a label, respond and close are events the tracker records anyway, and there is nothing to synchronise. Second, most items are small: content corrections, wording, missing links, configuration---exactly the class an agent resolves end to end with a reviewable pull request. Third, maintainer attention is the binding constraint, not code-review capacity: the agent shifts effort from \emph{producing} a first response to \emph{judging} one, and judging a two-line diff takes a minute. Under these conditions the human role reduces to merge, accept and confirm, and the loop closes without anyone owning a queue.

\subsection{Where it breaks}
The same conditions delimit the design. Items that need a decision the project has not made (scope, policy, roadmap) go down the proposal path and wait for a human exactly as before; the agent makes the waiting \emph{visible} (a proposal with a timestamp) but not shorter. Requesters without a tracker account cannot accept a proposal, so the reviewer acts for both---acceptable for a small site, not for a product with many stakeholders. A project with several teams needs routing beyond a label. And an agent that opens pull requests is only as safe as the review culture around it: \ffrs{} makes the review checkpoint structural, but a maintainer who merges without reading has removed it.

\subsection{Removing the database was the right cut}
The first design had a relational database, an outbox and a metrics view. Replacing them with the tracker as system of record removed one external service, one secret, one schedule and roughly a third of the code, and it made every metric recomputable by anyone with read access. What it cost was buffering: capture is now synchronous with the tracker, and an outage yields an honest \texttt{502} and a client-side retry rather than a queued item. For a small project this is the better trade---the tracker's availability is already the project's availability---and it keeps the whole system explainable in one diagram.

\subsection{What the build taught us about agents in the loop}
Two production-only findings shaped the agent's operating envelope. The first agent run produced a correct fix within eight minutes and then failed on permissions: it could read the organisation's repositories but not write, and it reported this precisely instead of retrying. The second run succeeded but spent most of its time on test infrastructure that did not match the sandbox. Both point to the same lesson: an agentic Respond stage needs a budget and an escape hatch (``state what you could not run and ship the pull request'') more than it needs cleverness, and the surrounding tooling---installable browsers, hermetic tests---matters as much as the prompt. Once both were in place, a real request was resolved by an agent pull request and a one-minute human review.

\subsection{Cost}
At the reference site's volume the whole system runs within a cloud free tier and the agent within a subscription; the marginal cost of an item is the reviewer's minute. We report actual figures in Sec.~\ref{sec:eval} once the measurement window closes.
